\subsection{Natural Gas Data}
The results of the natural gas data analysis highlight the need for using varying cycle lengths in the development of surrogate models for underground hydrogen storage. The variations in cycle lengths across Europe are likely due to variations in the reservoir properties, market demands, operational models and varying seasons. The six month idle phases observed in the Netherlands are due to the site being at maximum capacity \citep{AGSI_data}. The withdrawal cycle median being lowest in Italy is potentially due to having a shorter winter period. This is typically  when gas is produced to meet increased market demand. Another explanation for the low median withdrawal cycle length is the reservoir properties of the  Cornegliano gas storage site. This is the only site in Italy used in this study and it has the lowest total gas capacity of all sites \citep{AGSI_data}. This may result in rapid depletion of the storage and shorter withdrawal cycle lengths. A final explanation is the site's ``5+2'' operational model during the withdrawal season in winter. This model means the site withdraws gas for five days during the week and injects gas for two days at the weekend \citep{Italy_UGS}. If injection outweighs the withdrawal amount, the withdrawal periods may appear much shorter. This operational model is designed to optimize storage operations. Seasonal variations across the different countries potentially explains the observation that the median injection cycle length is lowest in the United Kingdom. The summer months are generally when gas is injected. As a result, the United Kingdom having a shorter summer period would explain the shorter injection cycle lengths. 

The results and analysis highlight the significant variations in cycle lengths and the multiple impacting variables. Due to this, training a deep learning surrogate would allow for rapid assessment of varying cycle lengths. The cyclic nature of underground hydrogen storage means deep learning models can be trained to assess these uncertainties once and then used multiple times. Although not clearly distinguishable in the results, this rapid assessment of uncertain cycle lengths would be invaluable for hydrogen storage operators when outlier events occur. Outlier events such as the invasion of Ukraine and Covid-19, can have a large impact on energy demand and assessing these impacts could be crucial. 

\subsection{Prediction of Scalar Values}

The injection and withdrawal rate are important to predict. These are typically operator defined in natural gas storage and are a function of reservoir pressure and physical properties of the geological formation \citep{brkic2021underground}. In underground gas storage minimum and maximum operating pressures are defined. This is to ensure reservoir integrity and caprock tightness \citep{wlodek2026underground}.

In this study wellbore hole pressure limits were used to comply with this standard. As a result, the Eclipse input injection or withdrawal rate can differ compared to the Eclipse output rate. When the wellbore hole pressure in the eclipse simulation reaches the maximum or minimum pressures, the rates are automatically adjusted to ensure the limits are not breached. Therefore the model predicting the injection and withdrawal rates allows the adjusted rate to be predicted instantly.

In relation to use of the model by an operator of a hydrogen storage site, they can input an injection or withdrawal schedule and then receive the achievable injection or withdrawal rate. This allows hydrogen storage operators to forecast the amount of energy they can store and provide. 

The prediction of injection rates is mostly very accurate. The accuracy drops above 10 $\times 10{^3} \ (kg \cdot M / D)$. This is likely due to their being less data points in this range, to limit this the DenseLoss was used \citep{steininger2021density}. The ability of the model to predict withdrawal rate is significantly lower than injection rate predictions. The $R^2$ value is only slightly better than predicting the mean. An explanation for this is that wellbore hole pressure changes much faster during withdrawal than injection which the model may struggle to learn (Appendix \ref{Appendix:WBHP_box}). To try and limit the impact of this, the three wellbore hole pressure components were added; however, this had a limited impact. The hydrogen stored component was added to the scalar inputs to try address this issue. The model predicts the amount of hydrogen stored with high accuracy. This is unexpected due to the models weak ability to predict withdrawal rate. The most plausible explanation is the difference in scales. The hydrogen stored prediction will not be affected significantly if the model learns to decrease during a withdrawal phase by the mean amount as it has much larger values. Despite the values being normalised, the difference in scales will have an impact. As the hydrogen stored is predicted independently, it will learn the withdrawal rate implicitly and the effects of it in relation to its own scale.

A hydrogen storage operator using the model can accurately predict the rate they will be able to store hydrogen and how much will be stored. This is crucial for forward planning. Accurately forecasting hydrogen deliverability will not be possible.

\subsection{Spatial Prediction}

The pressure predictions are highly accurate. Comparing the observation of underprediction in the north-east corner to the horizontal permeability, an explanation can be hypothesised. The horizontal permeability is higher than the surrounding area and this contrast may not be learnt by the model (fig. \ref{fig:Horizontal_perm}). Overall, the pressure predictions, although errors appearing wide-spread, are highly accurate. 

The saturation errors are highest closest to the well. This is likely due to it being where saturation will be most variable in injection and withdrawal phases. The observation of underprediction at the base of the anticline is potentially due to it being the mean plume front. Due to the hydrogen being driven upwards, the base of the anticline will have lower saturation. As the plume front will differ for all sites, the model does not learn to capture plume front physics entirely. There is also a layer of relatively high overpredictions by the well. When comparing to the horizontal permeability it can be seen that there is an area of lower permeability. The model therefore is unable to completely determine the difference in permeability and the effects of this. 

The errors in both saturation prediction and pressure prediction could  be reduced by inputting the reservoir properties into the model. However, this is unnecessary due to the high accuracy and the increased compute capacity it would require. The model can accurately predict the hydrogen plume shape and pressure within the reservoir. For hydrogen storage operators this is needed to avoid leakage of hydrogen and ensure reservoir integrity.

\subsection{Autoregressive Rollout}

The results of the injection rate autoregressive rollout predictions potentially give us an understanding of the key input variables to the model. The observation that the error increases initially loosely aligns with observations seen in the hydrogen stored and pressure autoregressive rollout results. This implies that these three variables receive and react to similar input signals resulting in errors propagating through these variables. The interquartile range significantly increasing after month 90 potentially indicates the results have gone significantly outside the model's training parameter space. The withdrawal rate autoregressive rollout prediction errors show no clear link to the autoregressive rollout inputs. This is not unexpected due to the results of the teacher-forced prediction being poor. This highlights that the model potentially does not receive enough information from the current inputs to make accurate predictions. Cumulative errors in the autoregressive rollout inputs do not propagate through to the withdrawal rate autoregressive rollout predictions. 

The saturation autoregressive rollout prediction error does not follow a trend shown by other variables. This result implies it solely depends on the conditioning variables and saturation inputs.  The saturation error increases smoothly which implies the prediction is not impacted by multiple variables. Another likely explanation is the size of the plume. In the first few months, the plume will be relatively small and throughout the time period it increases. Therefore, the error may grow due to the plume extent increasing and predictions becoming less accurate further from the well.

Overall, the model predictions from autoregressive rollout have a low error rate up to 15 months for injection rate, hydrogen stored and pressure. This means that operators can accurately predict and assess the impacts of injection cycles for the maximum cycle length observed in the natural gas data from one initial set of inputs. However, operators should be cautious in using the model to detect saturation leakage and forecasting deliverability.


\subsection{Limitations and Directions for Future Work}

The primary limitation of this study is the model used to generate the dataset. With no existing commercial operations of hydrogen storage in depleted gas fields, the accuracy of the model is not able to be verified. In addition to this, the resolution of the grid cells being 50m may not fully capture the dynamics of the hydrogen plume. Although there is no existing research on the resolution required for underground hydrogen storage, using carbon capture and storage as an analogue, resolutions of one to two metres are required \citep{gegewen_nested}. 

Future work should primarily focus on improving the accuracy of predicting withdrawal rate, the key value for forecasting deliverability. One method to explore would be using a much finer temporal resolution. This would likely reduce the impact of large wellbore hole pressure changes, which is potentially causing the prediction to be inaccurate. Investigating the impacts of higher temporal resolution on autoregressive rollout predictions is also an interesting avenue to explore in future. A final, and most valuable, avenue for future research would be testing autoregressive predictions on samples starting at different stages in the simulation and outside of the model's training time range. Achieving high accuracy on this is a very strong justification for developing deep learning based surrogate models. This would allow hydrogen storage operators from one single input to predict multiple future timesteps. 